\noindent\textbf{The base and the family.}
Let $\Delta = \Delta(3,4,\infty)$ be the triangle group, acting on the upper half plane $\hb$; the quotient $\hb/\Delta$ is $\dbP^1$ minus a point, and adding the cusp gives $\dbP^1$ with orbifold points of orders $3$ and $4$ and the cusp. Let $V \cong \dbZ^4$ carry automorphisms $T_1, T_2$ of orders $3$ and $4$ with $T_0 := (T_1 T_2)^{-1}$ unipotent, $(T_0 - I)^2 = 0$: this is a representation $\Delta \to \mathrm{SL}(V)$, and we let it act on the dual lattice $\Lambda = V^*$ by $A_j = (T_j^{-1})^t$ (\S2). The period functions $\tau, \mu, \beta$ on $\hb$ are an equivariant period map for it (\S3), and the quotient is a family of complex 2-tori over the punctured orbifold curve: the analogue of a universal family of abelian surfaces over a modular curve, with this rank-four lattice representation in place of the symplectic one. The very general fibre is not algebraic. The very general fibre $F$ has $\mathrm{NS}(F) = \dbZ\eta$ with $\eta$ a single indefinite class, and the algebraic dimension of the completed threefold is $a(X) = 1$ (\S9).

\medskip
\noindent\textbf{The three special fibres.}
The three conjugacy classes of maximal finite and parabolic cyclic subgroups of $\Delta$ single out the three special points, and at these three points we take the three standard completions; each is a choice, and every $\ell_0$ and every admissible $v_j$ below occurs. At an elliptic point of order $m$ the local monodromy is a torus automorphism of order $m$, and the filling is Kodaira's logarithmic transform: $(\mathrm{disc} \times \mathrm{torus})/(\dbZ/m)$, the generator acting by a rotation of the disc and by the automorphism followed by a translation; the reduced fibre is of bielliptic type (\S5). At the cusp the monodromy is unipotent of rank two, and the filling $N_0$ is Mumford's toric degeneration; its central fibre $W$ is read off the prescribed $A_2$-triangulation of $\dbR^2$, the deck action being the one in which the unimodular map $B_0 \colon \bar\Lambda \to \Lambda_{\mathrm{tor}}$ induced by $M_0 - I = (T_0^{-1})^t - I$ on $\Lambda$ enters. Modulo the lattice the triangulation has one vertex, three edges and two triangles, i.e.\ $W$ is one toric surface with hexagonal moment polygon --- the degree-six del Pezzo $\mathrm{dP}_6$ --- glued to itself along opposite sides, with two triple points (\S4). Thus there is a single family, completed in the three standard ways at the three kinds of special point of an orbifold curve with a cusp.

\medskip
\noindent\textbf{The gluing and the fundamental group.}
Each filling is glued to the global family along a collar by a fibre-preserving map, and we reglue by translations by local sections (\S6, Prop.~6.3); among these, $\pi_1$ and $H_*$ depend only on three integers $\ell_0$ (at the cusp), $\ell_1, \ell_2$ (at the elliptic points) (\S7). Let $\gamma \in V$ be the linear coordinate on $\Lambda$ whose kernel $\ker\gamma = \langle \hat u, \hat w, \hat\delta\rangle$ is the span of all $(A_j - 1)\Lambda$; then the coinvariants $\Lambda/\ker\gamma \cong \dbZ$ are detected by $\gamma$ alone, and $\ell_j = \gamma(v_j)$ for the translation vectors $v_j$. The result is $\pi_1(X) \cong \dbZ/|12\ell_0 - 4\ell_1 - 3\ell_2|$ (\S7).\footnote{\S7 proves $\pi_1(X) \cong \dbZ/|p|$, $p = 12\ell_0 - 4\ell_1 - 3\ell_2$, and computes $H_*(X;\dbZ)$ for every admissible pair $(v_1,v_2)$; the value $|p|=1$ is one point of that family.} This formula admits a reading, which the body does not use. It is the formula $|a_1 a_2 e_0 + a_2 b_1 + a_1 b_2|$ of \cite{Orl72} for the order of $H_1$ of the Seifert fibred space over $S^2(a_1,a_2)$, evaluated at $e_0 = \ell_0$ and $b_j = -\ell_j$ --- the sign coming from the convention $s_j \circ g_j = e^{-2\pi i/m_j} s_j$ for the local sections --- so here $(a_1,a_2) = (3,4)$ with Seifert invariants $(3,-\ell_1)$, $(4,-\ell_2)$ and obstruction $\ell_0$; for $(\ell_0,\ell_1,\ell_2) = (0,1,-1)$ the value is $1$ and the space in question is $S^3$ with the circle action $(z,w) \mapsto (\lambda^3 z, \lambda^4 w)$, whose orbit space is $S^2(3,4)$ \cite{Orl72}. Heuristically the twist condition says: glue so that the coinvariant circle of the fibres traces out the $(3,4)$ Seifert fibration of the three-sphere over the base.

\medskip
\noindent\textbf{The homology and the Euler number.}
The remaining torus directions contribute nothing: $\hat w$ and $\hat\delta$ are vanishing cycles at the cusp, collapsed in the hexagonal fibre, and the rest is killed by monodromy, the coinvariants of $\Lambda$ having rank one. With the twists above, $X$ has the integral homology of $S^6$ and is diffeomorphic to $S^6$ (\S7, \S8). The Euler number localises at the cusp: $e(X) = e(W) = 2$ (\S7), every other fibre being a torus or a free torus quotient and so of Euler number $0$ --- one singular fibre carries the whole Euler characteristic of $S^6$, as the twelve nodal fibres carry the $12$ of a rational elliptic surface.

\medskip
\noindent\textbf{Why the argument of \cite{CDP20} does not apply.}
That argument requires a line bundle non-torsion on every fibre and propagates a vanishing statement through each singular fibre by way of its normalisation. Here $H^2(X;\dbZ) = 0$ and $\mathrm{Pic}(X) \cong \dbC$, so every line bundle is topologically trivial; at the non-normal toric fibre the differential of the fibration itself supplies the section which the passage to the normalisation assumes away, and $R^2 f_*(T_X \otimes L) \ne 0$ for all $L$. The count of $H_1$ in \cite[Lemma 4.2]{CDP20} (after \cite[Lemma 3.2]{CDP98}) assumes trivial monodromy, whereas the monodromy representation here is non-trivial; \S10 shows that this last point is repairable for $X$, while the failure of the reduction at the non-normal fibre is decisive.

\bigskip
\noindent\textbf{Setup.}
$V := \dbZ^4$, basis $(\gamma,u,w,\delta)$; $\Lambda := V^*$, dual $(\hat\gamma,\hat u,\hat w,\hat\delta)$; $T_1, T_2 \in \mathrm{Aut}(V)$, orders $3,4$ (columns: images), and $\Pi(z)$ ($\tau,\mu,\beta$ below):
\[
T_1 = \begin{pmatrix} 1 & 0 & -6 & 2 \\ 0 & -1 & 1 & 1 \\ 0 & -1 & 0 & 1 \\ 0 & 0 & 0 & 1 \end{pmatrix}, \qquad
T_2 = \begin{pmatrix} 1 & 6 & 0 & -3 \\ 0 & 0 & -1 & 1 \\ 0 & 1 & 0 & 0 \\ 0 & 0 & 0 & 1 \end{pmatrix}, \qquad
\Pi(z) := \begin{pmatrix} 6\mu & \tau & 1 & 0 \\ \beta & \mu & 0 & 1 \end{pmatrix}.
\]
$T_0 := (T_1 T_2)^{-1} = I + N$, $N^2 = 0$, $N\gamma = Nu = 0$, $Nw = -u$, $N\delta = \gamma$; on $\Lambda$: $A_j := (T_j^{-1})^t$, $M_0 := (T_0^{-1})^t$, $\Lambda_{\mathrm{tor}} := \langle \hat w, \hat\delta\rangle$, $\bar\Lambda := \Lambda/\Lambda_{\mathrm{tor}}$, $B_0 \colon \bar\Lambda \xrightarrow{\sim} \Lambda_{\mathrm{tor}}$, $\bar{\hat\gamma} \mapsto -\hat\delta$, $\bar{\hat u} \mapsto \hat w$; $A_1, A_2$ fix $\varepsilon := \hat\gamma + 2\hat u - 4\hat w$, $\varepsilon' := \hat\gamma + 3\hat u - 3\hat w$ resp. $B := \dbP^1 \ni t$, $t_c := 1/t$, $p_1,p_2,p_0 := 0,1,\infty$, $B^\circ := B \setminus \{p_0,p_1,p_2\}$, $m_1,m_2 := 3,4$; $\pi \colon \hb_z \to B \setminus \{p_0\}$ ($\hb_z, \hb$ upper half-plane) orbifold uniformisation ($m_j$ at $p_j$, cusp $p_0$), deck group $\Delta = \langle g_1,g_2 \mid g_1^3 = g_2^4 = 1\rangle$, $g_j z_j = z_j \in \pi^{-1}(p_j)$, $g_j'(z_j) = e^{-2\pi i/m_j}$, $g_0 := (g_1 g_2)^{-1}$ parabolic; $\rho_V \colon \Delta \to \mathrm{GL}(V)$, $g_j \mapsto T_j$, $M_g := \rho_V(g)^t$; $U_r \subset \pi^{-1}\{0 < |t_c| < r\}$ ($r$ small) $\langle g_0\rangle$-invariant cusp component; holomorphic $\tau \colon \hb_z \to \hb$, $\mu,\beta \colon \hb_z \to \dbC$ ($j$ modular invariant, $j(i) = 1728$):
\[
(\tau 1)\colon j(\tau) = 1728\, t\circ\pi, \qquad
(\tau 2)\colon \tau\circ g_1 = \frac{\tau - 1}{\tau}, \ \ \tau \circ g_2 = -\frac1\tau,
\]
\[
(\tau 3)\colon \tau(z_1) = e^{i\pi/3}, \ \ \tau(z_2) = i, \qquad
(\mu 1)\colon \mu\circ g_1 = \frac{\mu}{1-\tau}, \ \ \mu\circ g_2 = \mu\left(1 + \frac{\tau}{2}\right),
\]
\[
(\mu 2)\colon \mu \text{ bounded on } U_r, \qquad
(\beta 1)\colon \beta\circ g_1 = \beta + 2 - \frac{6(1-\mu)^2}{\tau}, \ \ \beta\circ g_2 = \beta - 3 - \frac{6\mu\tau}{2},
\]
\[
(\beta 2)\colon \beta + \tau \text{ bounded on } U_r, \qquad
(\beta 3)\colon \operatorname{Im}\beta - \frac{6(\operatorname{Im}\mu)^2}{\operatorname{Im}\tau} < 0 \text{ on } \hb_z.
\]
Such $\tau,\mu,\beta$ exist by Theorem 3.4; on $\dbC^2 \ni \zeta = (\zeta_1,\zeta_2)$, $\Pi(z)$ as above with columns $\Pi(z)\hat\gamma, \Pi(z)\hat u, \Pi(z)\hat w, \Pi(z)\hat\delta$: $\lambda\cdot(z,\zeta) := (z, \zeta + \Pi(z)\lambda)$, unique $R_g(z) \in \mathrm{GL}_2(\dbC)$: $\Pi(gz) = R_g(z)\Pi(z)M_g$, $\mathcal T := (\hb_z \times \dbC^2)/\Lambda$, $\Delta$ by $\tilde g(z,\zeta) := (gz, R_g(z)\zeta)$, $J := (\mathcal T|_{\hb_z \setminus \Delta\{z_1,z_2\}})/\Delta \to B^\circ$, monodromy $A_1, A_2, M_0$ on $H_1 = \Lambda$. On $U_r$, $e^{2\pi i \tau} = t_c u_1(t_c)$, $u_1$ unit at $0$; for $s := \tau - (2\pi i)^{-1}\log u_1$ (any branch; \S4) $\Pi = [sB_0 + C(t_c) \mid I]$ in bases $(\bar{\hat\gamma},\bar{\hat u})$, $(\hat w,\hat\delta)$, $C$ holomorphic at $0$. $N' := \dbZ^3 \ni (y,y_3)$, $y \in \dbZ^2 = \Lambda_{\mathrm{tor}}$ ($e_1 \leftrightarrow \hat w$, $e_2 \leftrightarrow \hat\delta$), cocharacter lattice of $(\dbC^*)^3 \ni (x_1,x_2,t_c) \subset Y$ smooth toric, fan $\mathfrak F$ cone over $A_2$-triangulation of $\dbR^2 \times \{1\}$ (vertices $\dbZ^2 \times \{1\}$, edges $\pm e_1, \pm e_2, \pm(e_1 - e_2)$), $x_k = e^{2\pi i \zeta_k}$, $t_c$ character of $(0,0,1) \in \mathrm{Hom}(N',\dbZ)$, holomorphic on $Y$. $\Lambda$ acts freely, properly discontinuously on $Y_\epsilon := \{|t_c| < \epsilon\}$, $\epsilon$ small (4.5), by $\Psi_{\bar\lambda} := (\exp 2\pi i\, C(t_c)\bar\lambda, 1)\cdot \Phi_{\bar\lambda}$, $\Phi_{\bar\lambda}$ from $(y,y_3) \mapsto (y + y_3 B_0\bar\lambda, y_3)$; $N_0 := Y_\epsilon/\Lambda$, $W := \{t_c = 0\}/\Lambda$, $\eta \colon \mathrm{dP}_6 \to W$ normalisation, $\mathrm{dP}_6$ del Pezzo of degree $6$. For $j=1,2$: $\Delta_j \ni z_j$ $g_j$-invariant disc, coordinate $s_j$, $s_j \circ g_j = e^{-2\pi i/m_j} s_j$, $s_j^{m_j} = t_j \circ \pi$, $t_j$ coordinate at $p_j$, $N_j := (\mathcal T|_{\Delta_j})/\langle (z,\zeta) \mapsto (g_j z, R_{g_j}(z)\zeta + \Pi(g_j z) v_j/m_j)\rangle$ ($N_1$, $N_2$), $v_1 := \varepsilon$, $v_2 := -\varepsilon'$ (free; 5.4). Glue to $J$, any branch: $N_j$ by $\zeta \mapsto \zeta + (2\pi i)^{-1}(\log s_j)\Pi(z)v_j$ (\S5), $(Y_\epsilon \cap (\dbC^*)^3)/\Lambda$ on $0 < |t_c| < \epsilon$ by $(z,\zeta) \mapsto (e^{2\pi i \zeta_1}, e^{2\pi i \zeta_2}, e^{2\pi i s(z)})$. $\ell_j := \gamma(v_j)$ ($\gamma \in V = \mathrm{Hom}(\Lambda,\dbZ)$): $\ell_1 = 1$, $\ell_2 = -1$, $\ell_0 := 0$. $X := (J \sqcup N_0 \sqcup N_1 \sqcup N_2)/\!\sim$ (both gluings), $f \colon X \to B$ induced.

\begin{theorem}
$X$ is a compact connected complex $3$-manifold and $f \colon X \to B$ is a surjective holomorphic map with connected fibres.
\begin{enumerate}[label=(\arabic*)]
\item Over $B^\circ$ the map $f$ is a proper submersion whose fibres are complex $2$-tori. The fibre $f^{-1}(p_0)$ is the divisor $W \subset N_0$, which is reduced, irreducible and has normal crossings. The normalisation $\eta$ identifies opposite sides of the hexagon of $(-1)$-curves of $\mathrm{dP}_6$, the double locus is three rational curves through two triple points, and $e(W) = 2$. For $j=1,2$ one has $f^*(p_j) = m_j S_j$, where $S_j := (f^{-1}(p_j))_{\mathrm{red}}$ is a smooth bielliptic surface in $N_j$, and $\mathcal O_X(S_j)|_{S_j}$ has order $m_j$.
\item The threefold is simply connected, since $\pi_1(X) \cong \dbZ/|12\ell_0 - 4\ell_1 - 3\ell_2| = 1$, and $H_*(X;\dbZ) \cong H_*(S^6;\dbZ)$; consequently $X$ is diffeomorphic to $S^6$ (\S7ff.).
\item The algebraic dimension is $a(X) = 1$ and $\mathcal M(X) = f^*\dbC(t)$: every meromorphic function on $X$ is constant along the fibres of $f$.
\item The canonical bundle is $K_X \cong f^*\mathcal O_B(-1) \otimes \mathcal O_X(2S_2)$ and is not torsion; $c_3(X) = 2$; and $\mathrm{Aut}^0(X) \cong \dbC^*$ acts along the fibres of $f$ with fixed locus a double curve of $W$.
\end{enumerate}
\end{theorem}

\begin{remark}[on \cite{CDP20}]
The Theorem contradicts \cite[Cor.\ 2.3]{CDP20} as published. For $L \in \mathrm{Pic}(X)$ set $A := (L^* \otimes K_X)|_W$: $\eta^*A$ is trivial ($H^2(X;\dbZ)=0$, $\mathrm{Pic}(\mathrm{dP}_6)$ discrete), and for $\theta$ trivialising it $d(t_c \circ f)|_W \otimes \theta$ vanishes on the double locus, so descends to $0 \ne \sigma \in \Omega^1_X|_W \otimes A$ dying in $\widetilde\Omega^1_W \otimes A$ ($\widetilde\Omega := \Omega/\mathrm{torsion}$). By Serre--Grothendieck duality on $W$ ($\omega_W \cong K_X|_W$) and base change $R^2 f_*(T_X \otimes L) \ne 0$ for all $L$ (\S10); so hypothesis (1) of \cite[Prop.\ 2.4, p.\ 680]{CDP20} holds for no $L$. The reduction of \cite[p.\ 692]{CDP20} to $H^0(\widetilde\Omega^1_S \otimes \cdot) = 0$ on reduced fibres $S$ is stated provided $L|_S$ is non-torsion, a proviso met at $W$ for general $L$, where the conclusion $H^0(\widetilde\Omega^1_W \otimes A) = 0$ does hold; but it does not suffice. Indeed, for general $L$, $A \not\cong \mathcal O_W$ forces $H^0(W, N_W^* \otimes A) = H^0(W,A) = 0$, so the conormal term vanishes identically, while $\sigma$ gives $H^0(W, \Omega^1_W|_W \otimes A) \ne 0$: this is where the reduction fails.
\end{remark}

\newpage
